\documentclass[12pt,a4paper]{article}
\usepackage[T1]{fontenc}\usepackage[utf8]{inputenc}\usepackage{times}
\usepackage[a4paper,top=2.5cm,bottom=2.5cm,left=2.5cm,right=2.5cm]{geometry}
\usepackage{graphicx}
\graphicspath{{figures/}}
\usepackage{booktabs}\usepackage{array}\usepackage{enumitem}\usepackage{tabularx}
\usepackage{caption}\usepackage{fancyhdr}\usepackage{parskip}\usepackage{float}
\usepackage{setspace}\onehalfspacing
\usepackage{amsmath}
\usepackage{xurl}\usepackage[hidelinks]{hyperref}
\captionsetup{font=small,labelfont=bf,labelsep=colon,
              justification=justified,singlelinecheck=false}
\setlist[itemize]{leftmargin=1.3em,itemsep=0pt,parsep=0pt,topsep=2pt}
\setlist[enumerate]{leftmargin=1.6em,itemsep=0pt,parsep=0pt,topsep=2pt}
\setlength{\parskip}{4pt plus 1pt}\setlength{\emergencystretch}{2em}
\renewcommand{\topfraction}{0.85}\renewcommand{\bottomfraction}{0.7}\renewcommand{\textfraction}{0.1}\renewcommand{\floatpagefraction}{0.75}
\pagestyle{fancy}\fancyhf{}
\fancyhead[L]{\footnotesize KOUAME Koffi Fid\`ele}
\fancyhead[C]{\footnotesize NEW YORK CITY AIRBNB ANALYTICS}
\fancyhead[R]{\footnotesize Data Analysis Internship}
\fancyfoot[C]{\thepage}\renewcommand{\headrulewidth}{0.4pt}
\setlength{\headheight}{15pt}
\newcommand{\fig}[3][0.92]{\begin{figure}[H]\centering\includegraphics[width=#1\linewidth]{#2}\caption{#3}\end{figure}}
\newcommand{\capture}[2]{\IfFileExists{figures/#1}{\begin{figure}[H]\centering\setlength{\fboxsep}{0pt}\setlength{\fboxrule}{0.4pt}%
\fbox{\includegraphics[width=\dimexpr\linewidth-0.8pt\relax]{#1}}\caption{#2}\end{figure}}{}}

\begin{document}

\begin{titlepage}\centering
\singlespacing
\vspace*{3cm}
{\large\scshape Data Analysis Internship}\\[0.4cm]
{\normalsize Task 12: Final Data Analysis Project}\\[2.0cm]
\rule{\linewidth}{0.5pt}\\[0.8cm]
{\LARGE\bfseries New York City Airbnb Market Analysis}\\[0.6cm]
{\large Where, What and at Which Price?\\ Listing Types, Areas, Prices, Demand and Hosts\\ of 48{,}895 Listings (2019)}\\[0.8cm]
\rule{\linewidth}{0.5pt}\\[1.2cm]
{\normalsize Data Understanding $\bullet$ Cleaning $\bullet$ Exploratory Analysis $\bullet$ Price Model
$\bullet$ Power BI Dashboard $\bullet$ Insights and Recommendations}\\[3.0cm]
{\large\itshape Prepared by}\\[0.5cm]
{\Large\bfseries KOUAME Koffi Fid\`ele}\\[0.6cm]
{\large 8 October 2026}
\vfill\end{titlepage}

\singlespacing
\tableofcontents\newpage
\onehalfspacing

% =====================================================================
\section{Project overview}

\textbf{Context.} Airbnb has become a large part of the accommodation offer of New York City.
Three audiences need to understand this market: \textbf{hosts}, who want to know where and at
which price to list; \textbf{Airbnb}, which follows supply and demand in each area; and
\textbf{the city}, which follows how short-term rentals spread across neighbourhoods.

\textbf{Objective.} This final project carries a complete analysis, from the raw file to an
interactive dashboard, on \textbf{one dataset only, the one supplied}: 48{,}895 Airbnb listings
of New York City in 2019. It answers the seven questions of the brief:
\begin{enumerate}[label=Q\arabic*.]
\item Which listing types are the most common?
\item Which areas have the most activity?
\item How do prices differ across areas?
\item Is there a relationship between price, ratings and reviews?
\item Which factors seem to affect price?
\item Which areas or listings perform best?
\item Are there noticeable patterns or trends?
\end{enumerate}

\textbf{Deliverables.} The notebook \texttt{NYC-Airbnb-Analytics.ipynb} reproduces every number and
figure of this report; the cleaned dataset and the cleaning log are in \texttt{data/processed/};
the result tables in \texttt{results/}; the Power BI dashboard, with its DAX measures, Power Query
code and theme, in \texttt{powerbi/}; the final presentation (20 slides) in \texttt{presentation/}.

\textbf{Summary of the results.}
\begin{itemize}
\item After cleaning, \textbf{48{,}645 listings} of \textbf{37{,}283 hosts} remain (99.49\% of the file).
A typical night costs \textbf{\$105} (median); the listings received \textbf{1.14 million reviews}.
\item \textbf{Entire homes (52\%) and private rooms (46\%)} make the market, at a median of \$160 and \$70.
\item \textbf{Manhattan and Brooklyn hold 85\% of the listings}; Brooklyn receives the most reviews.
\item \textbf{Prices fall with the distance from Manhattan}: \$149 in Manhattan, \$65 in the Bronx.
\item \textbf{Room type and location set the price}: an entire home costs +119\% and a Manhattan
listing +37\%, other things equal; the model explains 52\% of the variation.
\item \textbf{Price and reviews are almost unrelated} ($\rho = -0.05$); the file has no ratings.
\item \textbf{The busiest listings are cheap rooms near the airports in Queens.}
\item \textbf{40\% of the listings look inactive}, and \textbf{0.9\% of the hosts run 10\% of the listings}.
\end{itemize}

% =====================================================================
\section{Dataset information}

\subsection{Source and scope}

The file \texttt{AB\_NYC\_2019.csv}, supplied for the project and the only dataset used, describes \textbf{48{,}895 listings} offered in New York City, one row per listing,
with \textbf{16 columns}. The last review date of the file is 8 July 2019, which dates the snapshot.
The listings belong to \textbf{37{,}457 hosts} and \textbf{221 neighbourhoods} in the five boroughs.
A grey map of the city (\texttt{New\_York\_City\_.png}) is supplied with the data and is used as
the background of the maps.

\subsection{Columns, meaning and types}

\begin{table}[H]\centering\small
\caption{Columns of the file, with the type read by pandas and the missing values.}
\begin{tabularx}{\linewidth}{@{}lXlr@{}}
\toprule
Column & Meaning & Type & Missing\\\midrule
id & listing identifier & integer & 0\\
name & title of the listing & text & 16\\
host\_id & host identifier & integer & 0\\
host\_name & first name of the host & text & 21\\
neighbourhood\_group & borough (Manhattan, Brooklyn, Queens, Bronx, Staten Island) & text & 0\\
neighbourhood & neighbourhood (221) & text & 0\\
latitude, longitude & position of the listing & decimal & 0\\
room\_type & Entire home/apt, Private room, Shared room & text & 0\\
price & asked price per night (US\$) & integer & 0\\
minimum\_nights & minimum stay required & integer & 0\\
number\_of\_reviews & reviews received since the listing was created & integer & 0\\
last\_review & date of the last review & text (date) & 10{,}052\\
reviews\_per\_month & average reviews per month & decimal & 10{,}052\\
calculated\_host\_listings\_count & listings of the same host in the file & integer & 0\\
availability\_365 & days available for booking in the next 365 days & integer & 0\\
\bottomrule
\end{tabularx}
\end{table}

\subsection{What the data can and cannot answer}

The file describes the \textbf{offer} (where, what, at which price, under which conditions) and gives
a \textbf{proxy of demand}: reviews. It has \textbf{no rating column}, no booking, no revenue and no
size of the property (bedrooms, capacity). Question Q4 on ratings is therefore answered with reviews,
which measure demand, not satisfaction. Prices are \emph{asked} prices per night.

Every numeric column is strongly skewed: the median price is \$106 but the mean is \$153 and the
maximum \$10{,}000; the median minimum stay is 3 nights but the maximum 1{,}250. The analysis therefore
uses \textbf{medians} to describe a typical listing and rank-based tests.

% =====================================================================
\section{Data cleaning process}\label{sec:clean}

\subsection{Issues found}

\begin{table}[H]\centering\small
\caption{Data quality audit and treatment.}
\begin{tabularx}{\linewidth}{@{}Xrl@{}}
\toprule
Issue & Rows & Treatment\\\midrule
\texttt{last\_review} and \texttt{reviews\_per\_month} missing & 10{,}052 & structural: rate set to 0, flag \texttt{has\_reviews}\\
Listing name missing & 16 & \emph{Unnamed listing}\\
Host name missing & 21 & column dropped (personal data)\\
Spaces around listing names & 254 & trimmed\\
\texttt{last\_review} stored as text & 38{,}843 & converted to a date\\
Price = \$0 & 11 & removed\\
Price above \$1{,}000 per night (6 at \$9{,}999 or \$10{,}000) & 239 & removed from the analysis set\\
Minimum nights above 365 (up to 1{,}250) & 14 & capped at 365\\
Duplicates (on the id and on all other columns) & 0 & none found\\
Map file \texttt{.png} that is in fact a JPEG & 1 file & read by its content\\
\bottomrule
\end{tabularx}
\end{table}

The missing values are not errors: the review date and the review rate are missing \textbf{exactly}
for the 10{,}052 listings with 0 reviews, because no review date exists for them. Setting the rate to
0 is therefore exact, and the date stays empty. A listing cannot be free, so the 11 prices at \$0
are errors; the 239 prices above \$1{,}000 (0.49\%) include placeholders (\$9{,}999, \$10{,}000) and
luxury outliers that would distort averages. They stay in the raw file. Two values that look extreme
are \textbf{kept} because they are plausible: 36\% of the listings have 0 days of availability
(calendar closed or fully booked) and one listing has 58.5 reviews per month (a hotel-like room in
the Theater District). The 989 listings that share a title with another one are different listings
of the same host, not duplicates.

\subsection{Cleaning log and cleaned dataset}

\begin{table}[H]\centering\small
\caption{Cleaning log (\texttt{data/processed/cleaning\_log.csv}), also shown in the dashboard.}
\begin{tabularx}{\linewidth}{@{}rlXrr@{}}
\toprule
Step & Name & Action & Rows after & Median price\\\midrule
1 & Raw file & \texttt{AB\_NYC\_2019.csv} as supplied & 48{,}895 & \$106\\
2 & Duplicates & checked on the id and on all other columns: none & 48{,}895 & \$106\\
3 & Columns & \texttt{host\_name} dropped (16 columns $\to$ 15) & 48{,}895 & \$106\\
4 & Text & 254 names trimmed; 16 missing names filled & 48{,}895 & \$106\\
5 & Types & \texttt{last\_review} converted to a date & 48{,}895 & \$106\\
6 & Missing reviews & 10{,}052 never-reviewed listings: rate = 0 & 48{,}895 & \$106\\
7 & Price = 0 & 11 listings removed & 48{,}884 & \$106\\
8 & Price $>$ \$1{,}000 & 239 listings removed & 48{,}645 & \$105\\
9 & Minimum nights & 14 values capped at 365 & 48{,}645 & \$105\\
10 & Derived columns & bands, host size, recency, sort orders & 48{,}645 & \$105\\
\bottomrule
\end{tabularx}
\end{table}

The cleaned dataset (\texttt{airbnb\_nyc\_clean.csv}) has \textbf{48{,}645 listings (99.49\%)}
and \textbf{29 columns}: the 15 original columns kept and 14 derived columns used by the analysis and
the dashboard: a flag \emph{has reviews}, price band (7 bands), minimum-stay class (1-2, 3-6, 7-29, 30+ nights), availability band
(0, 1-90, 91-180, 181-270, 271-365 days), host size (1, 2-5, 6+ listings), year of the last review,
days since the last review, \emph{active in the last 12 months} (reviewed at most 365 days before
8 July 2019) and sort orders. The median price moves only from \$106 to \$105: the cleaning does
not change the picture of a typical listing.

% =====================================================================
\section{Exploratory analysis}

\subsection{Key performance indicators}

\begin{table}[H]\centering\small
\caption{KPIs of the cleaned dataset (also the KPI cards of the dashboard).}
\begin{tabularx}{\linewidth}{@{}lrX@{}}
\toprule
KPI & Value & Comment\\\midrule
Listings & 48{,}645 & in 221 neighbourhoods\\
Hosts & 37{,}283 & 1.3 listings per host on average\\
Median / average price per night & \$105 / \$141 & the mean is pulled up by expensive listings\\
Total reviews & 1{,}136{,}096 & demand proxy\\
Listings with at least one review & 79.6\% & \\
Reviewed in the last 12 months & 59.8\% & the clearly active share\\
Median availability & 44 days a year & mean 112 days\\
Listings of multi-listing hosts & 33.9\% & hosts with 2 listings or more\\
Most common room type & Entire home/apt & 51.8\%\\
\bottomrule
\end{tabularx}
\end{table}

\subsection{Which listing types are the most common? (Q1)}
\fig{fig01_room_types}{Listings by room type with their median price (left) and room-type mix in each borough (right).}
\textbf{Entire homes and apartments are the most common type (25{,}216 listings, 51.8\%)}, closely
followed by \textbf{private rooms (22{,}273, 45.8\%)}; shared rooms are marginal (1{,}156, 2.4\%).
The two main types are \textbf{two different products}: an entire home costs a median \$160 per
night, a private room \$70, a shared room \$45. The mix depends on the borough: entire homes dominate
in Manhattan (61\%), private rooms in Queens and the Bronx (60\%), where hosts rent a room in their
own home. Private rooms receive almost as many reviews as entire homes (47\% against 51\% of all
reviews) with fewer listings.

\subsection{Which areas have the most activity? (Q2)}
\fig{fig02_boroughs}{Supply (share of listings), demand (share of reviews) and reviews per listing by borough.}
\textbf{Manhattan and Brooklyn concentrate the activity}: 85\% of the listings and 83\% of the
reviews. Manhattan has the most listings (44.2\%), but \textbf{Brooklyn receives the most reviews}
(42.8\% against 39.9\%). The small boroughs are used more intensively: \textbf{Staten Island (31.1
reviews per listing), Queens (27.7) and the Bronx (26.0)} beat Manhattan (21.1), and their median
review rate (0.76 to 1.00 per month) is about three times that of Manhattan (0.29).
\fig{fig03_neighbourhoods}{The 12 neighbourhoods with the most listings (left) and the most reviews (right).}
The ten largest neighbourhoods hold 48\% of the listings. \textbf{Williamsburg} (3{,}910 listings)
and \textbf{Bedford-Stuyvesant} (3{,}700) lead the supply, and \textbf{Bedford-Stuyvesant is the most
reviewed neighbourhood} of the city (109{,}862 reviews, ahead of Williamsburg with 85{,}412).
\fig{fig04_map}{Map of the listings by borough (left) and by price per night (right, log colour scale), on the map supplied with the data.}
The maps show the listings spreading from Midtown and Lower Manhattan into north Brooklyn and western
Queens, and the highest prices in Lower and Midtown Manhattan.

\subsection{How do prices differ across areas? (Q3)}
\fig{fig05_price_distribution}{Distribution of the price per night (left, log scale) and share of listings by price band (right).}
The median night costs \textbf{\$105} and the mean \$141; half of the listings ask between \$69 and
\$175, 35\% ask \$50 to \$99 and only 8\% ask \$300 or more.
\fig{fig06_price_borough_room}{Price per night by borough and room type (box = middle 50\% of the listings, label = median, log scale).}
By borough, the median price is \textbf{\$149 in Manhattan}, \$90 in Brooklyn, \$75 in Queens and
Staten Island and \textbf{\$65 in the Bronx}. For an entire home: Manhattan \$190, Brooklyn \$145,
Queens \$120, Bronx and Staten Island \$100. Inside every borough an entire home costs about twice a
private room, and the boxes overlap: location explains only part of the price.
\fig{fig07_neighbourhood_prices}{Median price of an entire home by neighbourhood (30+ listings): the 10 most and the 10 least expensive.}
By neighbourhood, for an entire home, \textbf{Tribeca (\$305)}, the Flatiron District, NoHo and
SoHo are the most expensive; Midwood (\$88) and Queens Village (\$97) the least expensive: a ratio
of \textbf{3.5 to 1} for the same type of product.

\subsection{Is there a relationship between price, ratings and reviews? (Q4)}
\fig{fig08_price_reviews}{Price and number of reviews of every reviewed listing (left, hexbin, log scales), median review rate (middle) and share of never-reviewed listings (right) by price band.}
The file has \textbf{no rating column}, so the question is answered with reviews. \textbf{Price and
reviews are almost unrelated}: Spearman $\rho = -0.05$ with the number of reviews and with the review
rate (significant because the sample is large, negligible in size); inside each room type $\rho$ stays
between $-0.13$ and $0.00$. The \$50-99 band has the highest median review rate (0.46 per month); above
\$200 the rate falls to 0.17 for \$500-1{,}000, and the share of \textbf{never-reviewed} listings rises
from 18\% around \$50-150 to \textbf{37\% above \$500}. Many reviews do not justify a higher price,
and demand concentrates on affordable listings.

\subsection{Which factors seem to affect price? (Q5)}
\fig[0.86]{fig09_price_factors}{Effect of each factor on the price per night, other factors held constant (multiple regression of the log price, 95\% confidence intervals). Reference: private room, Brooklyn, 1-2 nights, 0 days available, host with 1 listing.}
A multiple regression of the log price measures each factor with the others held constant. It
explains \textbf{52\%} of the variation of the price (59\% with the 221 neighbourhoods instead of the
five boroughs).
\begin{itemize}
\item \textbf{Room type is the strongest factor}: alone it explains 41\% of the variation. An entire
home costs \textbf{+119\%} compared with a private room, a shared room $-33\%$.
\item \textbf{Location comes second}: the neighbourhood alone explains 29\%, the borough 14\%. Compared
with Brooklyn, a comparable listing costs \textbf{+37\% in Manhattan}, $-13\%$ in Queens and $-24\%$
in the Bronx and Staten Island.
\item \textbf{Minimum stay}: listings that require 30+ nights are 26\% cheaper per night (7-29 nights:
$-18\%$): monthly rentals with a lower nightly rate.
\item \textbf{Availability}: listings open most of the year are 13\% to 34\% more expensive than closed
ones, the pattern of full-time listings.
\item \textbf{Host size has no effect} once room type and location are taken into account, and more
reviews go with slightly lower prices ($-5\%$ per unit of log reviews).
\end{itemize}
The remaining 48\% depends on what the file does not contain: size, quality, amenities.

\subsection{Which areas or listings perform best? (Q6)}
\fig[0.86]{fig10_performance}{Median price and median review rate of the 98 neighbourhoods with 50+ listings (bubble = number of listings); dashed lines: city medians.}
Performance is measured by demand, the review rate. Among the 98 neighbourhoods with 50+ listings,
the highest median review rates are \textbf{East Elmhurst (3.76 reviews per month, next to LaGuardia
airport)} and \textbf{Springfield Gardens (2.53, next to JFK)}, then Queens Village, Mott Haven (Bronx),
Jamaica and Flushing, all at a median price of \$60 to \$80: \textbf{cheap and busy}. Large central
areas (Williamsburg, Midtown, the Upper West Side) have many listings but low median rates (0.13 to
0.30), and the expensive Manhattan neighbourhoods (Tribeca, NoHo, Flatiron) sit in the lower right
corner. At the listing level, 63 of the 100 listings with the highest review rate are in
\textbf{Queens} and 64 are \textbf{private rooms}, often titled after the airport (``5 mins from
JFK'', ``LGA airport''). The two Theater District rooms at the top (58.5 and 28 reviews per month)
look like hotel rooms sold on Airbnb. Two models work: \textbf{volume} (affordable rooms near transport)
and \textbf{value} (expensive central homes, fewer stays).

\subsection{Are there noticeable patterns or trends? (Q7)}
\fig{fig11_recency_availability}{Listings by month of their last review (left) and share of listings by availability in the next 365 days (right).}
\textbf{A large part of the supply is inactive.} Only \textbf{60\%} of the listings were reviewed in the
last 12 months and 20\% were never reviewed; the last review of 52\% of all listings falls in 2019, with
a sharp peak in June-July 2019 (the active listings). \textbf{36\% of the listings have 0 days of
availability}, while 20\% are open almost all year. Demand goes with availability and short stays:
the review rate rises with availability ($\rho = +0.30$) and falls with the minimum stay
($\rho = -0.25$).
\fig{fig12_hosts}{Share of hosts and of listings by host size (left) and minimum stay by host size (right).}
\textbf{There are two kinds of hosts.} \textbf{86\% of the hosts have a single listing} and hold 66\%
of the offer, with a median availability of only 6 days: occasional home sharing. At the other end,
\textbf{343 hosts (0.9\%) run 6 listings or more} and hold \textbf{10\% of the offer} (4{,}755 listings),
open a median of \textbf{310 days} a year, and \textbf{47\% of their listings require 30+ nights}:
professional operators of furnished apartments (the 10 largest hosts alone run 1{,}268 listings).
The file is one snapshot, so trends over time can only be read through the date of the last review.

\subsection{Statistical tests}

Prices and reviews are skewed, so the tests are rank-based. With 48{,}645 listings even weak effects
are significant; the size of the effect is what matters.

\begin{table}[H]\centering\small
\caption{Kruskal-Wallis tests on the price, Spearman correlations (95\% bootstrap intervals) and Mann-Whitney test.}
\begin{tabularx}{\linewidth}{@{}Xrl@{}}
\toprule
Test & Effect & Reading\\\midrule
Price by room type (Kruskal-Wallis) & $\varepsilon^2 = 0.46$ & large\\
Price by neighbourhood (221 groups) & $\varepsilon^2 = 0.30$ & large\\
Price by borough & $\varepsilon^2 = 0.14$ & moderate\\
Price by host size / by minimum stay & $\varepsilon^2 = 0.03$ / $0.02$ & small\\
Price -- number of reviews & $\rho = -0.05$ [$-0.06$, $-0.04$] & negligible\\
Price -- reviews per month & $\rho = -0.05$ [$-0.06$, $-0.05$] & negligible\\
Reviews per month -- availability & $\rho = +0.30$ [0.29, 0.31] & moderate\\
Reviews per month -- minimum nights & $\rho = -0.25$ [$-0.26$, $-0.24$] & moderate\\
Entire homes: multi-listing vs single hosts (Mann-Whitney) & medians \$175 vs \$155 & multi-listing hosts ask more\\
\bottomrule
\end{tabularx}
\end{table}

All price differences are significant ($p < 0.001$). The regression of Section~4.6 uses robust (HC3)
standard errors; all effects of room type, borough, minimum stay (except 3-6 nights) and availability
are significant, host size is not. These are associations, not causes.

% =====================================================================
\section{Important visualizations}

Each question has one visual chosen for what it must show.

\begin{table}[H]\centering\small
\caption{Choice of the visuals.}
\begin{tabularx}{\linewidth}{@{}llX@{}}
\toprule
Question & Visual & Why\\\midrule
Q1 listing types & bars + 100\% stacked bars & counts with the median price on each bar; the mix in each borough\\
Q2 active areas & side-by-side bars, map & supply and demand next to each other; the map shows where\\
Q3 prices by area & histogram, box plots, ranked bars & skewed prices need a log scale and the spread, not only an average\\
Q4 price and reviews & hexbin + two bar charts & 38{,}726 points would hide each other; bands show the trend, one scale per chart\\
Q5 price factors & dot plot with intervals & one effect per factor, other factors held constant\\
Q6 best performers & bubble chart & price, demand and size of each neighbourhood at once\\
Q7 patterns & time line, bars, stacked bars & recency, availability and host profiles\\
\bottomrule
\end{tabularx}
\end{table}

The \textbf{five most important visuals} are: the room-type mix (Figure~1), because it shows that the
market is two products; the supply and demand by borough (Figure~2), because Brooklyn leads demand;
the price by borough and room type (Figure~6); the price factors (Figure~9), because it ranks what
sets the price; and the demand-price bubble chart (Figure~10), because it shows where listings
perform best. The figures follow one standard: one fixed colour per borough (a categorical palette
checked for colour blindness, the same hues in the notebook, the dashboard and the presentation), a
fixed colour per room type (entire home coral, private room cyan, shared room gold), labelled values
placed so that no text overlaps another, one scale per chart, medians rather than means, and the base
of each chart in its title.

% =====================================================================
\section{Dashboard}\label{sec:dash}

The Power BI dashboard (\nolinkurl{powerbi/NYC_Airbnb_Dashboard.pbix}) is built on the two cleaned
tables of Section~\ref{sec:clean}, the listings and the cleaning log, so every figure it shows is a
figure of this report. Its model holds 44 DAX measures: the indicators of the analysis (listings,
hosts, median and average price, reviews per listing, share of listings reviewed in the last 12
months, share of multi-listing hosts...), the medians restricted to groups large enough to be
compared (30 listings for a price, 50 for a review rate), the city median used as a reference, and
the counts of the cleaning log. Its seven pages follow the questions of the brief
(Table~\ref{tab:pages}).

\begin{table}[H]\centering\small
\caption{\label{tab:pages}Pages of the dashboard.}
\begin{tabularx}{\linewidth}{@{}lX@{}}
\toprule
Page & Content\\\midrule
1. Home & photo of the New York skyline; eight KPI cards (listings, hosts, median and average price, reviews, reviews per listing, share active in the last 12 months, neighbourhoods); word cloud of the neighbourhoods by number of listings; three picture tiles that open the analysis pages; a scrolling band of the largest neighbourhoods\\
2. Map \& Boroughs & map of the 48{,}645 listings by borough; listings by borough; listing types; median price by borough (Q1, Q2)\\
3. Prices & median price by borough and room type; listings by price band; gauge of the median against the city median; the ten most expensive neighbourhoods; median price by minimum stay (Q3, Q5)\\
4. Demand & demand against price by neighbourhood; the ten busiest listings; median reviews per month by borough and by price band; listings by year of last review (Q2, Q4, Q6)\\
5. Hosts & shares of hosts and of listings by host size; minimum stay by host size; listings per host and hosts with 6 or more listings; days open in the year; share of multi-listing hosts by borough (Q7)\\
6. Insights & the nine findings and the recommendations\\
7. Data Pipeline & raw, removed and clean listings, share kept, cleaning steps; the steps from the supplied file to the model and the cleaning log\\
\bottomrule
\end{tabularx}
\end{table}

\textbf{Interactivity.} Five dropdown slicers (borough, neighbourhood, room type, price band, minimum
stay) run along the top of the four analysis pages and keep their selection from one page to the next.
Clicking a bar, a slice or a point filters the other visuals of the page, and the cards, the gauge and
the top-ten visuals follow every selection: the gauge, for example, compares the median of the selected
listings with the median of the whole city. A page navigator in the header links the seven pages, and
the picture tiles of the Home page open the Map, Prices and Demand pages. Tooltips add a second figure to each bar (for example the share of listings and the median price of a borough), and every visual carries an alternative text for screen readers.

\textbf{Design.} A dark navy theme with a single blue accent: a header band with a soft blue glow, the
logo, the title and the page navigator; framed visuals with their title; KPI cards with rounded borders.
Single measures are blue; a borough and a room type keep the same hue on every page as in the figures of
this report, lightened for the dark background. The Home page uses the photo of the New York skyline
supplied with the project, and the map plots each listing on the map of New York City supplied with the
dataset. The icons are Microsoft's 3D Fluent Emoji (MIT licence).
Cards, axes and tables use business labels (``Median price'', ``Share of listings'') rather than the technical
names of the measures; value axes are removed where the values are written on the bars. In the model, the
44 measures are grouped in display folders (key figures, prices, demand, hosts and availability, shares,
data quality) with a description, and the technical sort columns are hidden.

\capture{dashboard/p1_home.png}{\label{fig:p1}Home: 48{,}645 listings, 37{,}283 hosts, a median of \$105 and an average of \$141.31 per night, 1{,}136{,}096 reviews (23.4 per listing), 59.8\% of listings active in the last 12 months, 221 neighbourhoods; Williamsburg and Bedford-Stuyvesant are the largest neighbourhoods.}
\capture{dashboard/p2_map_boroughs.png}{\label{fig:p2}Map \& Boroughs: Manhattan (21{,}488) and Brooklyn (20{,}041) hold 85\% of the listings; 51.8\% entire homes and 45.8\% private rooms; median of \$149 in Manhattan and \$65 in the Bronx.}
\capture{dashboard/p3_prices.png}{\label{fig:p3}Prices: an entire home costs \$190 in Manhattan and \$100 in the Bronx; one listing in three asks \$50 to \$99; without a filter the gauge sits on the city median of \$105; Tribeca is the most expensive neighbourhood (\$282, all room types).}
\capture{dashboard/p4_demand.png}{\label{fig:p4}Demand: the cheap neighbourhoods of Queens get the most reviews a month; demand is highest in Staten Island (1.00) and lowest in Manhattan (0.29), and falls above \$100 a night; 25{,}148 listings were last reviewed in 2019.}
\capture{dashboard/p5_hosts.png}{\label{fig:p5}Hosts: 86.2\% of hosts have one listing and hold 66.1\% of the listings; the 343 hosts with 6 or more hold 9.8\%, 47\% of them with stays of 30 nights or more; 35.9\% of calendars closed.}
\capture{dashboard/p6_insights.png}{\label{fig:p6}Insights: the nine findings and the recommendations for hosts, Airbnb and the city.}
\capture{dashboard/p7_data_pipeline.png}{\label{fig:p7}Data Pipeline: 48{,}895 raw listings, 250 removed, 48{,}645 kept (99.49\%), ten cleaning steps.}

% =====================================================================
\section{Key findings}

\begin{enumerate}
\item \textbf{Two products}: entire homes (52\%, median \$160) and private rooms (46\%, \$70).
\item \textbf{Two leading boroughs}: Manhattan and Brooklyn hold 85\% of the listings; Brooklyn
receives the most reviews (42.8\%) and Bedford-Stuyvesant is the most reviewed neighbourhood.
\item \textbf{Intensive outer boroughs}: Staten Island 31 and Queens 28 reviews per listing, Manhattan 21.
\item \textbf{A price gradient}: median \$149 in Manhattan, \$65 in the Bronx; Tribeca \$305 and Midwood
\$88 for an entire home.
\item \textbf{What sets the price}: room type (+119\%), location (+37\% in Manhattan vs Brooklyn),
minimum stay ($-26\%$ for 30+ nights); host size has no effect.
\item \textbf{Price and reviews are unrelated} ($\rho = -0.05$); 37\% of listings above \$500 were never reviewed.
\item \textbf{The busiest areas are near the airports}: East Elmhurst 3.76 and Springfield Gardens 2.53
reviews per month, at \$60-80.
\item \textbf{40\% of the listings look inactive}; 36\% have a closed calendar.
\item \textbf{A small professional segment}: 0.9\% of hosts, 10\% of listings, open 310 days a year,
half with 30+ night stays.
\end{enumerate}

% =====================================================================
\section{Business insights}\label{sec:ins}

\begin{table}[H]\centering\small
\caption{Insights: what was discovered, what it means and why it is important.}
\begin{tabularx}{\linewidth}{@{}>{\raggedright\arraybackslash}p{3.9cm}>{\raggedright\arraybackslash}X>{\raggedright\arraybackslash}X@{}}
\toprule
What we discovered & What it means & Why it is important\\\midrule
Entire homes and private rooms make 98\% of the offer, at \$160 and \$70 & two products with two kinds of guests & prices and competitors must be compared within a room type\\
Brooklyn has fewer listings than Manhattan but more reviews & demand has moved beyond Manhattan & Brooklyn is as important as Manhattan for growth\\
Outer boroughs have more reviews per listing & supply is short where demand per listing is high & new hosts there face less competition\\
Price falls with the distance from Manhattan & location sets the price level & a city-wide average is a wrong benchmark\\
Room type and location explain most of the price; host size does not & hosts control the price through what and where they offer & pricing advice can be built on few variables\\
Price and reviews are unrelated; expensive listings are more often never reviewed & reviews follow affordability, not price & a high price does not bring more stays\\
Cheap rooms near the airports are the busiest listings & convenience and price win short stays & a clear, under-served niche\\
40\% of the listings were not reviewed in 12 months & headline counts over-state the active supply & search results and market size are inflated\\
0.9\% of hosts run 10\% of the listings, open all year & a professional segment behaves differently & the city can target checks without burdening home sharers\\
\bottomrule
\end{tabularx}
\end{table}

% =====================================================================
\section{Recommendations}

\textbf{For hosts}
\begin{enumerate}
\item \textbf{Price against the right benchmark}: the median of the same room type in the same
neighbourhood (Prices page), not the city average; a price far above it needs a visible reason.
\item \textbf{Consider the outer boroughs and transport hubs}: rooms near JFK and LaGuardia and in
Queens are booked far more often, at a median of \$60-80.
\item \textbf{Open the calendar and shorten the minimum stay} if the goal is short stays: demand rises
with availability ($\rho = +0.30$) and falls with the minimum stay ($\rho = -0.25$).
\end{enumerate}
\textbf{For Airbnb}
\begin{enumerate}[resume]
\item \textbf{Activate dormant supply}: 20\% of the listings were never reviewed and 36\% have a closed
calendar; reminders, pricing tips or delisting would make search results more reliable.
\item \textbf{Recruit hosts where demand per listing is high} (Queens, the Bronx, Staten Island) rather
than in saturated central areas.
\item \textbf{Segment the offer by guest type}: affordable rooms for short, convenience-driven trips;
entire homes in Manhattan for leisure and families; monthly rentals for longer stays.
\end{enumerate}
\textbf{For the city}
\begin{enumerate}[resume]
\item \textbf{Monitor the professional segment}: host size, availability and minimum stay identify the
operators of many entire homes (343 hosts with 4{,}755 listings) and allow targeted checks of the rules
on short-term rentals.
\end{enumerate}
\textbf{For the next analysis}
\begin{enumerate}[resume]
\item \textbf{Add ratings, bookings and property size} to move from demand proxies to satisfaction,
occupancy and revenue.
\end{enumerate}

\subsection*{Limitations}
\begin{itemize}
\item No ratings: Q4 is answered with reviews, which measure demand, not satisfaction.
\item No bookings or revenue: prices are asked prices, and not every guest leaves a review.
\item No property size or quality, which is why the price model stops at $R^2 = 0.52$.
\item One snapshot: time patterns come from the date of the last review only.
\item Listings above \$1{,}000 per night (0.5\%) are excluded: the luxury segment is not covered.
\end{itemize}

% =====================================================================
\section{Conclusion}

From a raw file of 48{,}895 listings, this project produced a clean and documented dataset (48{,}645
listings, 10 logged steps), answered the seven questions of the brief with figures and tests, and
delivered an interactive Power BI dashboard. The New York Airbnb market of 2019 is a market of
\textbf{two products} (entire homes and private rooms), \textbf{concentrated in Manhattan and
Brooklyn}, where \textbf{the room type and the location set the price}. Demand, measured by reviews,
does not follow the price: it is strongest for \textbf{affordable rooms in the outer boroughs and near
the airports}. A large part of the listed supply is inactive, and a small group of professional hosts
runs one listing in ten. These findings give hosts a pricing benchmark, Airbnb a map of where supply
is short or dormant, and the city a way to identify professional operators. Adding ratings, bookings
and property size would turn this demand analysis into an analysis of satisfaction and revenue.

\section*{Reproducibility}

Every number, table and figure of this report is produced by the notebook
\nolinkurl{NYC-Airbnb-Analytics.ipynb} from the files of \nolinkurl{data/raw/}. The notebook writes
the cleaned dataset and the cleaning log (\nolinkurl{data/processed/}), the result tables
(\nolinkurl{results/}) and the figures (\nolinkurl{figures/}). The dashboard
\nolinkurl{powerbi/NYC_Airbnb_Dashboard.pbix} reads \nolinkurl{data/processed/}. The final presentation is
\nolinkurl{presentation/NYC_Airbnb_Final_Presentation.pptx}. The project is
published at \url{https://github.com/koffifidelek59-collab/nyc-airbnb-analytics}.

\section*{References}
{\small\begin{enumerate}[itemsep=0pt]
\item Data Analysis Internship, Task 12. \emph{AB\_NYC\_2019.csv} and \emph{New\_York\_City\_.png}, dataset supplied for the final project.
\item Gebru, T. et al. (2021). Datasheets for Datasets. \emph{Communications of the ACM}, 64(12), 86--92.
\item Kruskal, W. H., \& Wallis, W. A. (1952). Use of ranks in one-criterion variance analysis. \emph{Journal of the American Statistical Association}, 47(260), 583--621.
\item Seabold, S., \& Perktold, J. (2010). statsmodels: Econometric and statistical modeling with Python. \emph{Proceedings of the 9th Python in Science Conference}.
\item Microsoft. Power BI documentation. \url{https://learn.microsoft.com/power-bi/}
\item Microsoft. \emph{Fluent Emoji} (3D icons of the dashboard), MIT License. \url{https://github.com/microsoft/fluentui-emoji}
\end{enumerate}}

\enlargethispage{4\baselineskip}
\vspace{0.1cm}
\noindent\hfill\begin{minipage}{7cm}\begin{flushright}\singlespacing
\textbf{KOUAME Koffi Fid\`ele}\\
Data Analysis Internship\\
\href{mailto:koffifidelek59@gmail.com}{koffifidelek59@gmail.com}
\end{flushright}\end{minipage}

\end{document}
